\cleardoublepage
\pdfbookmark[1]{Declarações}{declarations}
\chapter*{Declarações}

\noindent\textit{Esta tese é submetida ao concurso para o cargo de
Professor Titular na Universidade Federal de Santa Catarina (UFSC),
Departamento de Engenharia de Computação, campus Araranguá. Toda a investigação
foi conduzida de forma independente pelo candidato, como demonstração de
liderança acadêmica e maturidade investigativa exigidas neste nível da carreira
docente.}

\bigskip

\section*{Financiamento}

Esta pesquisa foi realizada com recursos institucionais da Universidade Federal
de Santa Catarina e com suporte de infraestrutura computacional do Laboratório
de Bioinformática e Engenharia de Sistemas Biológicos. As fontes de financiamento
não tiveram papel no delineamento do estudo, na coleta ou análise de dados, na
decisão de publicar nem na preparação do manuscrito.

\section*{Disponibilidade de Dados e Código}

Todos os dados e o código-fonte que sustentam os resultados desta tese
estão abertamente disponíveis e são reproduzíveis.

\begin{sloppypar}
A implementação de software SHYpn (Redes de Petri Hierárquicas de Sinais)
encontra-se disponível em \url{https://github.com/simao-eugenio/shypn},
licenciada sob os termos da MIT License. O modelo de esporulação de
\textit{Bacillus subtilis} e as redes de exemplo dos Capítulos~4 e~5
encontram-se em \texttt{workspace/validation/}. Todos os parâmetros cinéticos
foram recuperados da base BRENDA via API documentada; os dados de vias
metabólicas provêm da base KEGG via REST API; as fórmulas elementares foram
obtidas da ChEBI via servidor FTP. As instruções de instalação constam do
arquivo \texttt{README.md} no repositório.
\end{sloppypar}

\section*{Interesses Conflitantes}

O autor declara não possuir interesses financeiros ou não financeiros
conflitantes. Este trabalho foi realizado como investigação acadêmica
independente, sem filiação comercial ou patrocínio que pudesse ser percebido
como influência sobre a objetividade dos resultados.

\section*{Aprovação Ética}

Esta tese envolve modelagem computacional e simulação de sistemas biológicos.
Não foram utilizados sujeitos humanos, experimentos em animais nem dados
clínicos. Portanto, aprovação ética não foi requerida.

\section*{Contribuição do Autor}

Esta tese representa investigação original independente conduzida pelo autor.
Todas as contribuições teóricas (formalismo da 13-tupla, semântica de consumo de
sinal, teoria de preempção hierárquica, independência fraca), o desenvolvimento
do software SHYpn (mais de 15.000 linhas em Python, interface GTK, motores
híbridos de simulação), a validação quantitativa (reprodução da assimetria gradual/abrupto de
\textcite{Fujita2005}: ${\approx}\,52\%$ vs.\ ${\approx}\,5\%$, separatriz
de $\sigma^H$ emergente a $\theta_{\sigma H} \approx 1{,}60\,\mu$M, sem
ajuste de parâmetro) e a redação dos capítulos foram realizados de forma
autônoma pelo candidato.
Colaboradores externos, reconhecidos na seção de agradecimentos, ofereceram
especialidade em biologia e casos de teste. Curadores das bases de dados KEGG,
BRENDA e ChEBI viabilizaram a validação por meio de seus recursos mantidos com
rigor.

\section*{Publicações Decorrentes desta Tese}

\textbf{Artigos em periódicos:}
\begin{enumerate}
    \item \textit{[A completar: submissões a periódicos]}
\end{enumerate}

\textbf{Trabalhos em conferências:}
\begin{enumerate}
    \item \textit{[A completar: apresentações em congressos]}
\end{enumerate}

\noindent Todas as publicações estão devidamente reconhecidas onde material
pertinente aparece nesta tese.

\section*{Materiais de Terceiros}

Esta tese inclui figuras e dados das seguintes fontes, utilizados com
permissão ou sob licença adequada: diagramas de vias KEGG (licença acadêmica
KEGG); estruturas químicas ChEBI (Creative Commons Attribution 4.0); dados
enzimáticos BRENDA (termos de uso acadêmico); bibliotecas Python, NumPy e SciPy
(licenças BSD); bibliotecas GTK (LGPL). Todos os materiais de terceiros estão
devidamente citados no texto e atendem aos respectivos termos de licenciamento.

\bigskip
\bigskip

\noindent Assinatura: \rule{6cm}{0.4pt}

\noindent Data: \rule{4cm}{0.4pt}
